\documentclass[12pt]{article}
\usepackage{amsmath, amssymb, amsthm}
\usepackage{geometry}
\geometry{margin=2.5cm}

% Simple manual environments - no amsthm numbering involved
\newenvironment{satz}[1]{%
  \par\vspace{6pt}\noindent\textbf{Theorem #1.}\itshape\space
}{%
  \par\vspace{6pt}
}

\newenvironment{definition}[1]{%
  \par\vspace{6pt}\noindent\textbf{Definition #1.}\itshape\space
}{%
  \par\vspace{6pt}
}

\newenvironment{bemerkung}[1]{%
  \par\vspace{6pt}\noindent\textbf{Remark #1.}\itshape\space
}{%
  \par\vspace{6pt}
}

\renewcommand{\qedsymbol}{$\blacksquare$}
\setcounter{secnumdepth}{0}

\title{5.3 Taylor's Formula}
\date{}

\begin{document}
\maketitle

Let $I \subset \mathbb{R}$ be a perfect interval, $f \in C^n(I, \mathbb{R})$ and $x_0 \in \mathring{I}$.

\textbf{Goal:} Approximate $f$ ``as well as possible'' by a polynomial of the form
\[
    p(x) = \sum_{k=0}^{n} c_k(x-x_0)^k.
\]

\textbf{Question:} How do the coefficients $c_k$ depend on $f$?

\textbf{Idea:} Choose $c_k$ so that $p$ agrees with $f$ in as many derivatives as possible at $x_0$. It follows:
\begin{align*}
    f(x_0) = p(x_0) \quad &\text{d.h.} \quad c_0 = f(x_0), \\
    f'(x_0) = p'(x_0) \quad &\text{d.h.} \quad c_1 = f'(x_0), \\
    &\vdots \quad \vdots \quad \vdots \\
    f^{(k)}(x_0) = p^{(k)}(x_0) \quad &\text{d.h.} \quad c_k = \frac{f^{(k)}(x_0)}{k!}.
\end{align*}

Subsequent analysis applies itself unproblematically to the following general situation:

\begin{enumerate}
    \item $I$ is a convex, perfect subset of $\mathbb{K}$; hereby a subset $M \subset X$ ($X$ a normed vector space) is called convex if for all $x, y \in M$ the connecting line segment
    \[
        \{(1-t)x + ty : 0 \leq t \leq 1\} \subset M.
    \]

    \item $f$ is a mapping with values in a normed vector space $(Y, \|\cdot\|_Y)$.
\end{enumerate}

\begin{definition}{5.3.1}
Let $M \subset \mathbb{K}$ be convex and perfect and let $(Y, \|\cdot\|_Y)$ be a normed VR.

\begin{enumerate}
    \item Let $f \in C^n(M, Y)$, $n \in \mathbb{N}_0$. Then we call
    \[
        T_n(f, x_0)(x) := \sum_{k=0}^{n} \frac{1}{k!}f^{(k)}(x_0)(x-x_0)^k
    \]
    the $n$-th Taylor polynomial of $f$ with expansion point $x_0$.

    \item Let $f \in C^\infty(M, Y)$. Then we call
    \[
        T(f, x_0) = \sum_{k=0}^{\infty} \frac{1}{k!}f^{(k)}(x_0)(x-x_0)^k
    \]
    the Taylor series of $f$ with expansion point $x_0$.
\end{enumerate}
\end{definition}

% ...existing code...

\begin{satz}{5.3.2 (Taylor's Formula)}
\begin{enumerate}
    \item Let $I \subset \mathbb{R}$ be a perfect interval and $f \in C^n(I,\mathbb{R})$, with $f^{(n)}$ differentiable on $\mathring I$. Then for every $x,x_0\in I$ there exists $\xi\in\mathring I$ such that
    \[
        f(x)=\sum_{k=0}^{n}\frac{1}{k!}f^{(k)}(x_0)(x-x_0)^k
        +\frac{1}{(n+1)!}f^{(n+1)}(\xi)(x-x_0)^{n+1}.
    \]

    \item Let $M\subset\mathbb K$ be convex and perfect. Let $f\in C^n(M,\mathbb K^m)$ and let $f^{(n)}$ be differentiable on $\mathring M$. Then for every $x,x_0\in M$ there exists $\Theta\in\left]0,1\right[$ such that
    \[
        \left\|f(x)-\sum_{k=0}^{n}\frac{1}{k!}f^{(k)}(x_0)(x-x_0)^k\right\|_2
        \le
        \frac{1}{(n+1)!}\left\|f^{(n+1)}\!\big(x_0+\Theta(x-x_0)\big)(x-x_0)^{n+1}\right\|_2.
    \]
\end{enumerate}
\end{satz}

\begin{proof}
\begin{enumerate}
\item Let $x,x_0\in I$. W.l.o.g.\ $x\neq x_0$. Set
\[
\rho:=\big[f(x)-T_n(f,x_0)(x)\big]\frac{(n+1)!}{(x-x_0)^{n+1}},
\]
and
\[
\phi(t):=
f(x)-\sum_{k=0}^{n}\frac{1}{k!}f^{(k)}\!\big(x_0+t(x-x_0)\big)(x-x_0)^k(1-t)^k
-\rho(1-t)^{n+1}\frac{(x-x_0)^{n+1}}{(n+1)!}.
\]
Then $\phi(0)=0$ and $\phi(1)=0$. Also
\begin{align*}
\phi'(t)
&=
-\sum_{k=0}^{n}\frac{1}{k!}f^{(k+1)}\!\big(x_0+t(x-x_0)\big)(x-x_0)^{k+1}(1-t)^k \\
&\quad
+\sum_{k=0}^{n}\frac{k}{k!}f^{(k)}\!\big(x_0+t(x-x_0)\big)(x-x_0)^k(1-t)^{k-1}
+\rho(1-t)^n\frac{(x-x_0)^{n+1}}{n!} \\
&=
-\frac{1}{n!}f^{(n+1)}\!\big(x_0+t(x-x_0)\big)(x-x_0)^{n+1}(1-t)^n
+\rho(1-t)^n\frac{(x-x_0)^{n+1}}{n!}.
\end{align*}
By Rolle's theorem, $\exists\,\Theta\in\left]0,1\right[$ with $\phi'(\Theta)=0$, hence
\[
\rho=f^{(n+1)}\!\big(x_0+\Theta(x-x_0)\big).
\]
Substituting this into the definition of $\rho$ gives (1), with
\[
\xi:=x_0+\Theta(x-x_0)\in\mathring I.
\]

\item Let $x,x_0\in M$, $x\neq x_0$, and let $\rho$ be as in part (1). Note $\rho\in\mathbb K^m$. Define $v\in\mathbb K^m$ by
\[
v_k:=
\begin{cases}
\frac{1}{\sqrt m}, & \text{if }\rho=0,\\[2mm]
\frac{\rho_k}{\|\rho\|_2}, & \text{if }\rho\neq 0,\ 1\le k\le m.
\end{cases}
\]
Then
\[
\sum_{k=1}^m \rho_k v_k=\|\rho\|_2,\qquad \|v\|_2=1.
\]
Set
\[
\psi(t):=\sum_{l=1}^{m}v_l\left\{
f_l(x)-\sum_{k=0}^{n}\frac{1}{k!}f_l^{(k)}\!\big(x_0+t(x-x_0)\big)(1-t)^k(x-x_0)^k
-\rho_l\frac{(1-t)^{n+1}}{(n+1)!}(x-x_0)^{n+1}
\right\}.
\]
As in (1): $\psi(0)=\psi(1)=0$, and
\[
\psi'(t)=\sum_{l=1}^{m}v_l\{\rho_l-f_l^{(n+1)}(x_0+t(x-x_0))\}\frac{1}{n!}(x-x_0)^{n+1}(1-t)^n.
\]
Rolle implies $\exists\,\Theta\in\left]0,1\right[$ with $\psi'(\Theta)=0$. Therefore
\[
\|\rho\|_2=\sum_{l=1}^{m}\rho_l v_l
=\sum_{l=1}^{m}f_l^{(n+1)}\!\big(x_0+\Theta(x-x_0)\big)v_l
\le
\left\|f^{(n+1)}\!\big(x_0+\Theta(x-x_0)\big)\right\|_2\|v\|_2,
\]
hence
\[
\|\rho\|_2\le \left\|f^{(n+1)}\!\big(x_0+\Theta(x-x_0)\big)\right\|_2.
\]
With the definition of $\rho$, this yields
\[
\|f(x)-T_n(f,x_0)(x)\|_2
\le
\frac{1}{(n+1)!}|x-x_0|^{n+1}
\left\|f^{(n+1)}\!\big(x_0+\Theta(x-x_0)\big)\right\|_2.
\]
\end{enumerate}
\end{proof}

% ...existing code...

\begin{bemerkung}{5.3.3}
Let $I \subset \mathbb{R}$ be a perfect interval, $f \in C^n(I,\mathbb{R})$, and $x_0 \in \mathring I$. Assume
\[
f'(x_0)=f''(x_0)=f^{(3)}(x_0)=\cdots=f^{(n-1)}(x_0)=0
\quad\text{and}\quad
f^{(n)}(x_0)\neq 0.
\]
Then:
\begin{itemize}
    \item $n$ odd $\Longrightarrow$ $x_0$ is not a local extremum.
    \item $n$ even $\Longrightarrow$ $x_0$ is a local extremum, more precisely:
    \[
    f^{(n)}(x_0)>0 \Longrightarrow f(x_0)\ \text{local minimum},
    \qquad
    f^{(n)}(x_0)<0 \Longrightarrow f(x_0)\ \text{local maximum}.
    \]
\end{itemize}

\begin{proof}
Assume $f^{(n)}(x_0)>0$ (otherwise consider $-f$).
Since $f\in C^n(I,\mathbb{R})$, there exists $\delta>0$ with
\[
\left]x_0-\delta,\ x_0+\delta\right[\subset I
\quad\text{and}\quad
f^{(n)}(x)>0
\quad
\forall x\in\left]x_0-\delta,\ x_0+\delta\right[.
\]
By Theorem 5.3.2, for every $x\in\left]x_0-\delta,\ x_0+\delta\right[$ there exists
$\xi\in\left]x_0-\delta,\ x_0+\delta\right[$ such that
\[
f(x)=f(x_0)+\frac{f^{(n)}(\xi)}{n!}(x-x_0)^n.
\]
Hence:
\[
n\ \text{odd} \Longrightarrow
\forall x\in\left]x_0,\ x_0+\delta\right[:\, f(x)>f(x_0)
\ \text{and}\
\forall x\in\left]x_0-\delta,\ x_0\right[:\, f(x)<f(x_0),
\]
\[
n\ \text{even} \Longrightarrow
\forall x\in\left]x_0-\delta,\ x_0+\delta\right[:\, f(x)>f(x_0).
\]
\end{proof}
\end{bemerkung}

% ...existing code...

\begin{bemerkung}{5.3.4}
\begin{enumerate}
    \item The Taylor series of $f$ is a power series and converges in its radius of convergence. For $x \neq x_0$ in the interval of convergence:
    \[
        f(x) = T(f,x_0)(x)
    \]
    if and only if
    \[
        \lim_{n\to\infty}\left|f(x)-T_n(f,x_0)(x)\right|=0.
    \]

    \item For the function $f:\mathbb{R}\to\mathbb{R}$ with
    \[
        f(x)=
        \begin{cases}
            0, & x\le 0,\\
            e^{-1/x}, & x>0,
        \end{cases}
    \]
    Example 5.1.30 gives: $T(f,0)\equiv 0$.

    The function $T(f,0)(x)$ is defined for all $x\in\mathbb{R}$ (radius of convergence $\infty$), but for $x>0$ it does not agree with $f$ at any point.
\end{enumerate}
\end{bemerkung}

\begin{bemerkung}{5.3.5}
\begin{enumerate}
    \item Theorem 5.3.2(2) also holds for any other norm in $\mathbb{K}^m$, and also for functions with values in an arbitrary normed vector space $(Y,\|\cdot\|_Y)$.

    \item Assume the hypotheses of Theorem 5.3.2(2) with $n=0$. Then:
    \[
        \|f(x)-f(y)\|_Y
        \le
        \left(\max_{\xi\in\overline{xy}}\|f'(\xi)\|_Y\right)\,|x-y|.
    \]
\end{enumerate}
\end{bemerkung}

\end{document}